\subsection{Native audit and measured scope}

The driver \code{compressed\_word\_research/anchored\_replay.py} compares
against \code{b41fc4960968}, loading that package in an isolated namespace.
It pins 446 current source and fixture files before and after each run.
Among pre-existing runtime files only the compressed dispatcher changes;
the source-anchored checker is a new module. Producer discovery and the
literal checker are byte-for-byte unchanged.

All 1,132 maintained tests pass in 219.025 seconds. A separate audit covers
92 diagrams: the 19-case survivor/control corpus, filtered random closed
braids and five stabilized circles. Default, projection and forest discovery
return identical certificates in both packages. Every positive is replayed
by both the literal and compressed verifiers in both packages, giving 768
source replays. The new checker handles 74 blocks, all in projection mode;
the default and forest modes do not activate it in this audit. There is no
change in recognition coverage. The final pinned audit takes 8.988 seconds;
that duration is not a benchmark.

Three abstract infinite-cyclic presentation families additionally compare
exact retained words through run-length representations. These fixtures are
algebraic capacity tests, not diagram-derived recognition claims. A balanced
32-generator, 128-bit fixture takes five rounds: final ordinary nodes fall
from 8,170 to 3,341 and replay work from 119,154 to 45,156. A 31-round star
has 385 final nodes in both engines, but work falls from 86,179 to 27,944.
A mixed projection/forest chain with 16 generators and 512-bit local powers
has 22,662 versus 22,635 final nodes and 968,191 versus 315,907 replay work.
The largest reconstructed cumulative exponents have respectively 641, 129
and 7,681 bits. Exact words, slot positions and rank-one endpoints agree.

The three timing runs execute serially after tests, audit and first code
publication. Each uses five shuffled rounds, both engines and an A/A copy
of each engine. All 620 measured calls and 124 warm-ups complete; failed
calls would have been retained. The tables report median times and medians
of within-round ratios, which need not equal ratios of the median times.
Raw samples and deduplicated proofs are retained in the three
\code{data/anchored-replay-*.json} benchmark records; their table generator
is \code{data/anchored\_replay\_tables.py}.

\paragraph{Supplied algebraic schedules.}
Each timed call constructs a fresh grammar, authenticates every move of the
supplied schedule and verifies the final rank-one condition. Discovery and
exact output comparisons are outside the timer. The four cases give 80
measured calls and 16 warm-ups in 19.049 seconds.

\begin{center}\small
\begin{tabular}{lrrrrr}
\toprule
Case & Old ms & New ms & Old/new & Old A/A & New A/A \\
\midrule
star-4-0 & 0.342 & 0.283 & 1.215 & 1.017 & 1.003 \\
balanced-32-64 & 55.790 & 28.266 & 2.010 & 0.999 & 0.959 \\
star-32-128 & 79.009 & 39.927 & 1.979 & 1.010 & 0.996 \\
chain-16-512-mixed & 907.480 & 308.764 & 2.986 & 0.986 & 0.989 \\
\bottomrule
\end{tabular}
\end{center}


The three binary-power families give paired speed ratios of about $2.01$,
$1.98$ and $2.99$. Their complete endpoints agree, so these are not ratios
against an incomplete baseline. Their supplied schedules and abstract source
presentations nevertheless prevent interpreting them as full recognition
speedups.

\paragraph{Complete supplied certificates from diagrams.}
Each timed call validates a fresh PD, independently recovers its presentation,
and checks the entire supplied certificate through its endpoint. Certificate
discovery is excluded. The eight cases give 160 measured calls and 32 warm-ups
in 4.390 seconds. The number in each name is the stabilized circle's strand
count; it has one fewer crossings.

\begin{center}\small
\begin{tabular}{lrrrrr}
\toprule
Case & Old ms & New ms & Old/new & Old A/A & New A/A \\
\midrule
circle-5-projection & 0.343 & 0.345 & 1.012 & 0.958 & 1.071 \\
circle-5-forest & 0.324 & 0.303 & 1.073 & 1.054 & 0.960 \\
circle-17-projection & 2.213 & 2.100 & 1.069 & 1.001 & 1.024 \\
circle-17-forest & 1.733 & 1.752 & 0.981 & 1.001 & 1.014 \\
circle-65-projection & 21.290 & 19.167 & 1.116 & 0.995 & 1.000 \\
circle-65-forest & 14.043 & 14.087 & 0.988 & 1.001 & 0.998 \\
circle-129-projection & 75.447 & 68.209 & 1.123 & 1.003 & 0.977 \\
circle-129-forest & 47.267 & 47.237 & 0.998 & 1.004 & 0.981 \\
\bottomrule
\end{tabular}
\end{center}


The 129-strand projection certificate changes from 75.447 to 68.209 ms,
with paired old/new ratio $1.123$. Charged work falls from 171,431 to 153,791;
ordinary nodes remain 1,017 and maximum word-length bit size remains three.
The new route checks 63 blocks comprising 126 moves, visits 378 nodes for
lazy length queries, and uses only one-bit cumulative exponents. Both
engines make 126 LCP calls and the same 318 frontier queries visiting 698
nodes. This source-bound example gains through reduced repeated maintenance,
not through enormous hidden powers. The corresponding forest proofs contain
one forest batch and retain the direct checker; their ratios are effectively
one. The sub-millisecond five-strand cases are noisy, including A/A variation,
and do not support a reliable improvement claim.

\paragraph{Whole recognition, including discovery.}
Both engines enable primitive projections and compressed search on the
19-case corpus. Calls include fresh PD validation, discovery and mandatory
independent source replay, with the same 20-million group-work, ten-second
group, twelve-second total and 50,000-object limits. All 380 measured calls
and 76 warm-ups complete in 34.843 seconds.

\begin{center}\small
\begin{tabular}{lrrrrr}
\toprule
Case & Old ms & New ms & Old/new & Old A/A & New A/A \\
\midrule
survivor-00 & 15.430 & 15.481 & 1.012 & 0.995 & 0.997 \\
survivor-01 & 13.185 & 13.159 & 1.002 & 1.003 & 1.003 \\
survivor-02 & 14.496 & 14.192 & 1.023 & 0.986 & 1.005 \\
survivor-03 & 10.318 & 10.225 & 0.996 & 1.001 & 0.997 \\
mirror-03 & 23.029 & 22.907 & 1.007 & 0.993 & 0.990 \\
survivor-04 & 8.868 & 8.924 & 0.979 & 0.885 & 1.023 \\
survivor-05 & 17.787 & 17.895 & 0.994 & 0.993 & 1.005 \\
survivor-06 & 7.064 & 7.092 & 0.993 & 0.996 & 0.999 \\
survivor-07 & 19.071 & 18.637 & 1.018 & 1.008 & 1.023 \\
survivor-08 & 7.145 & 7.346 & 1.022 & 1.011 & 0.970 \\
mirror-08 & 19.440 & 19.207 & 1.009 & 0.998 & 0.997 \\
survivor-09 & 7.142 & 7.149 & 0.997 & 1.006 & 0.915 \\
survivor-10 & 6.102 & 6.063 & 1.013 & 1.015 & 0.997 \\
survivor-11 & 6.614 & 6.688 & 1.000 & 1.006 & 1.006 \\
gordian & 1240.018 & 1250.434 & 1.007 & 1.013 & 0.988 \\
trefoil & 0.056 & 0.056 & 0.993 & 1.026 & 1.009 \\
figure\_eight & 0.062 & 0.061 & 1.013 & 0.993 & 0.993 \\
conway & 1.015 & 0.992 & 1.014 & 1.024 & 1.002 \\
kinoshita\_terasaka & 1.049 & 1.017 & 1.003 & 1.023 & 1.002 \\
\bottomrule
\end{tabular}
\end{center}


Paired old/new ratios range from $0.979$ to $1.023$, with no broad improvement.
Old A/A ratios range from $0.885$ to $1.026$ and new A/A ratios from $0.915$
to $1.023$. The matching direct-discovery audit finds no eligible consecutive
monomial blocks on these 19 sources. This corpus therefore supplies a
compatibility and overhead control, not evidence for a whole-recognizer
speedup. Producer-side lazy discovery remains a separate integration task.
The native checker improvement proves neither a bound on general source
resets nor a general quasipolynomial unknot algorithm.
